\documentclass[10pt,a4paper]{amsart}
\usepackage[margin=19mm]{geometry}
\usepackage{amsmath,amssymb,mathtools}
\usepackage[scr=rsfs,cal=euler]{mathalfa}
\usepackage{xfrac}
\usepackage[unicode,hidelinks,bookmarks=false]{hyperref}
\newcommand{\ie}{i.e.\ }
\newcommand{\cf}{cf.\ }
\newcommand{\resp}{respectively\ }
\newcommand{\Subsection}{\subsection}
\providecommand{\nobreakdash}{\nobreak\hbox{-}}
\setlength{\emergencystretch}{2em}
\multlinegap0pt
\usepackage{mathtools}
\usepackage{xcolor}
\usepackage{tikz}
\usepackage{algorithm}\usepackage{algpseudocode}
\usepackage[normalem]{ulem}
\newtheorem{prop}{Proposition}
\newcommand{\poly}{\Omega}
\newcommand{\sizeA}{N}
\newcommand{\umax}{\mu_{1}}
\newcommand{\umin}{\mu_{\sizeA}}
\title{Block {Alpha-}Circulant Preconditioners for All-at-Once Diffusion-Based Covariance Operators}
\author{Jemima M. Tabeart, Selime G\"urol, John W. Pearson, Anthony T. Weaver}
\date{Source: arXiv:2506.03947v4 (CC BY 4.0)}
\begin{document}
\maketitle

\noindent\emph{Source and licence.} Block {Alpha-}Circulant Preconditioners for All-at-Once Diffusion-Based Covariance Operators. Version arXiv:2506.03947v4 (CC BY 4.0), distributed by arXiv under the Creative Commons Attribution 4.0 International licence, http://creativecommons.org/licenses/by/4.0/, as recorded in the arXiv licence metadata for this exact version and shown on the abstract page https://arxiv.org/abs/2506.03947v4. Full licence notice: \texttt{licenses/arxiv-2506.03947-CC-BY-4.0.txt}.\par\smallskip

\noindent\textit{Editorial scope.} Counted unit: the article's prop prop:errorbound (source0.tex lines 470--474). The article's complete original proof of it is delivered with it (source0.tex lines 476--489, one \texttt{proof} environment of 14 lines). No mathematics is written, repaired or re-derived: the statement and the proof are the article's own lines, in the article's own notation, and no retained line is substituted or re-flowed.  Retained uncounted as the source-local context the proof needs: lines 332--334; lines 411--413; lines 417--421.  Nothing was omitted from any retained span, so the package carries no omission record.  No retained line contains a citation, so the wrapper carries no bibliography and no bibliography entry.  No retained line contains a figure environment, an image-inclusion command or drawing code, so no image is delivered and the package declares no asset dependency.  Editorial formatting only, in addition: an AMS \texttt{amsart} preamble that restates the article's own declarations and macros used by the retained lines, namely the environments \texttt{prop} on separate counters, together with the standard packages those lines need.  The retained environments print in this document as Proposition 1 in order of appearance, whereas the article prints the same items with the numbering of its own full document.  The complete native source of the article as distributed in the arXiv e-print for this exact version is bundled unchanged as \texttt{sources/179300/source0.tex}.
		\begin{equation}\label{eq:B definition}
			B_j = {A} - \lambda_j I_\sizeA, \quad {j = 1, \ldots, \ell}.
		\end{equation}

		\begin{equation}\label{ep_error}
			\frac{\| {e}_p \|_2}{\| {e}_0 \|_2} \le \|  \poly_p(B_j) \|_2.
		\end{equation}

		\begin{align}
			\nonumber
			\|  \poly_p(B_j) \|_2 = \|  \poly_p(\Xi_j) \|_2 ={}& \text{spec} ( \poly_p(\Xi_j) ) \\
			\label{eq:norm_bound} ={}& \max_{k = 1, \ldots, \sizeA} |  \poly_p(\mu_k - \Re(\lambda_j) - i\; \Im(\lambda_j))  | \le \max_{{z} \in S_j} |  \poly_p( {z})|,
		\end{align}

		\begin{prop}\label{prop:errorbound}
			Let $B_j$ be defined as in \eqref{eq:B definition}. 
			Then CI converges when applied to $B_j$, with  the error at the $p$-th iteration being bounded above by \begin{equation*}
				\frac{\| {e}_p \|_2}{\| {e}_0 \|_2} \le {\left| T_p \left(\frac{\xi_j^{(1)}+\xi_j^{(N)}}{\xi_j^{(1)}-\xi_j^{(N)}}\right) \right|}^{-1}.\end{equation*}
		\end{prop}

		\begin{proof}
			The CI method takes $\poly_p$ in \eqref{ep_error} to be the shifted-and-scaled Chebyshev polynomial, i.e.
			\begin{equation*}
				{\poly_p(\mu) = \frac{T_p \left( \frac{\xi_j^{(1)} + \xi_j^{(N)} - 2 \mu}{\xi_j^{(1)} - \xi_j^{(N)}} \right)}{T_p \left( \frac{\xi_j^{(1)} + \xi_j^{(N)}}{\xi_j^{(1)} - \xi_j^{(N)}} \right)} =}  \frac{T_p \left(\frac{\umax+\umin-2\lambda_j-2\mu}{\umax-\umin}\right)}{T_p \left(\frac{\umax+\umin-2\lambda_j}{\umax-\umin}\right)},
			\end{equation*}
			where $\xi_j^{(1)}$ and $\xi_j^{(N)}$ are the eigenvalues of $B_j$ with largest and smallest real part, respectively, and $\mu \in \mathbb{C}$.
			Hence,
			\begin{align*}
				\max_{\mu \in [{\xi_j^{(\sizeA)},\xi_j^{(1)}}]} |  \poly_p( \mu) | ={}& \max_{\mu \in [{\xi_j^{(\sizeA)},\xi_j^{(1)}}]} \frac{\left| T_p \left(\frac{\umax+\umin-2\lambda_j-2\mu}{\umax-\umin}\right) \right|}{\left| T_p \left(\frac{\umax+\umin-2\lambda_j}{\umax-\umin}\right) \right|} \\
				={}& \frac{1}{\left|T_p \left( \frac{\xi_j^{(1)} + \xi_j^{(N)}}{\xi_j^{(1)} - \xi_j^{(N)}}  \right)\right|},
			\end{align*}
			by setting the argument {$\mu = \umin- \lambda_j$} to obtain the last line. {Using this relation, together with \eqref{ep_error} and \eqref{eq:norm_bound},
				we obtain the desired result.}
		\end{proof}

\end{document}
